\documentclass[10pt]{article}
\usepackage[utf8]{inputenc}
\usepackage[vietnamese]{babel}
\usepackage{amsmath,amssymb,amsfonts,amsthm}
\usepackage{geometry}
\geometry{
  paperwidth=210mm,
  paperheight=297mm,
  left=18mm,
  right=18mm,
  top=18mm,
  bottom=20mm,
  headheight=14pt
}
\usepackage{xcolor}
\usepackage{graphicx}
\usepackage{array}
\usepackage{colortbl}
\usepackage{fancyhdr}
\usepackage{newpxtext,newpxmath}

% Định nghĩa bảng màu theo sách Stewart Calculus
\definecolor{stewartcyan}{RGB}{0, 138, 179}
\definecolor{stewartred}{RGB}{204, 37, 41}
\definecolor{tableheader}{RGB}{232, 242, 246}
\definecolor{footergray}{RGB}{100, 100, 100}

% Cấu hình header & footer
\pagestyle{fancy}
\fancyhf{}
\fancyhead[L]{\textsf{\textbf{222}}\qquad \textsf{\textbf{CHƯƠNG 3}}\quad \textsf{Các quy tắc tính đạo hàm}}
\fancyhead[R]{}
\renewcommand{\headrulewidth}{0pt}
\fancyfoot[C]{%
  \fontsize{6.5pt}{8pt}\selectfont\color{footergray}%
  Bản dịch tiếng Việt dành cho mục đích học tập. Tác quyền tác phẩm gốc thuộc về © Cengage Learning.
}

\setlength{\parindent}{0pt}
\setlength{\parskip}{5pt}
\setlength{\fboxsep}{4.5pt}
\setlength{\fboxrule}{0.8pt}

\begin{document}

\noindent
% CỘT LỀ TRÁI (0.27\textwidth)
\begin{minipage}[t]{0.27\textwidth}
  \vspace{0pt}
  % Đồ thị Hình 4
  \centering
  \includegraphics[width=\linewidth]{images/p0257_fig1.png}
  
  \vspace{4pt}
  {\small\textbf{\textsf{HÌNH 4}}}
  
  \vspace{14pt}
  % Bảng giá trị
  {\footnotesize
  \renewcommand{\arraystretch}{1.22}
  \setlength{\tabcolsep}{4pt}
  \begin{tabular}{|l|l|}
    \hline
    \rowcolor{tableheader}
    \multicolumn{1}{|c|}{$x$} & \multicolumn{1}{c|}{$(1 + x)^{1/x}$} \\
    \hline
    $0{,}1$        & $2{,}59374246$ \\
    $0{,}01$       & $2{,}70481383$ \\
    $0{,}001$      & $2{,}71692393$ \\
    $0{,}0001$     & $2{,}71814593$ \\
    $0{,}00001$    & $2{,}71826824$ \\
    $0{,}000001$   & $2{,}71828047$ \\
    $0{,}0000001$  & $2{,}71828169$ \\
    $0{,}00000001$ & $2{,}71828181$ \\
    \hline
  \end{tabular}
  }
\end{minipage}%
\hfill
% CỘT NỘI DUNG CHÍNH (0.69\textwidth)
\begin{minipage}[t]{0.69\textwidth}
  \vspace{0pt}
  Vì $f'(1) = 1$, ta có
  \[
    \lim_{x \to 0} \ln(1 + x)^{1/x} = 1
  \]
  Khi đó, theo Định lý 2.5.8 và tính liên tục của hàm số mũ, ta có
  \[
    e = e^1 = e^{\lim_{x \to 0} \ln(1 + x)^{1/x}} = \lim_{x \to 0} e^{\ln(1 + x)^{1/x}} = \lim_{x \to 0} (1 + x)^{1/x}
  \]

  \vspace{2pt}
  \noindent
  \makebox[0pt][l]{%
    \fcolorbox{stewartred}{white}{\bfseries\sffamily\color{stewartred}\;5\;}%
  }%
  \hfill
  \fcolorbox{stewartred}{white}{%
    \hspace{1.6em}$\displaystyle e = \lim_{x \to 0} (1 + x)^{1/x}$\hspace{1.6em}%
  }%
  \hfill\null

  \vspace{5pt}
  Công thức 5 được minh họa bằng đồ thị của hàm số $y = (1 + x)^{1/x}$ trong Hình 4 và bảng giá trị ứng với các giá trị nhỏ của $x$. Điều này minh họa cho sự kiện là, chính xác đến bảy chữ số thập phân,
  \[
    e \approx 2{,}7182818
  \]
  Nếu ta đặt $n = 1/x$ trong Công thức 5, thì $n \to \infty$ khi $x \to 0^+$ và do đó một dạng biểu thức khác cho $e$ là

  \vspace{2pt}
  \noindent
  \makebox[0pt][l]{%
    \fcolorbox{stewartred}{white}{\bfseries\sffamily\color{stewartred}\;6\;}%
  }%
  \hfill
  \fcolorbox{stewartred}{white}{%
    \hspace{1.6em}$\displaystyle e = \lim_{n \to \infty} \left(1 + \frac{1}{n}\right)^{\!n}$\hspace{1.6em}%
  }%
  \hfill\null

  \vspace{12pt}
  % Mục Đạo hàm của các hàm số lượng giác ngược
  \noindent
  {\color{stewartcyan}\rule{1.1ex}{1.1ex}}\hspace{0.5em}{\large\textbf{Đạo hàm của các hàm số lượng giác ngược}}\par

  \vspace{3pt}
  Các hàm số lượng giác ngược đã được ôn tập trong Mục 1.5. Chúng ta đã thảo luận về tính liên tục của chúng trong Mục 2.5 và các đường tiệm cận của chúng trong Mục 2.6. Ở đây chúng ta sử dụng phương pháp đạo hàm ẩn để tìm đạo hàm của chúng. Ở phần đầu của mục này, ta đã nhận thấy rằng nếu $f$ là một hàm số đơn ánh và khả vi, thì hàm ngược $f^{-1}$ của nó cũng khả vi (ngoại trừ những điểm mà tiếp tuyến có phương thẳng đứng). Vì các hàm số lượng giác---với các tập xác định thu hẹp được sử dụng để định nghĩa hàm ngược---là đơn ánh và khả vi, nên suy ra các hàm số lượng giác ngược cũng khả vi.

  Nhắc lại định nghĩa của hàm arcsin:
  \[
    y = \sin^{-1} x \qquad \text{nghĩa là} \qquad \sin y = x \quad \text{và} \quad -\frac{\pi}{2} \le y \le \frac{\pi}{2}
  \]
  Lấy đạo hàm hai vế của $\sin y = x$ theo biến $x$ bằng đạo hàm ẩn, ta nhận được
  \[
    \cos y \, \frac{dy}{dx} = 1 \qquad \text{hay} \qquad \frac{dy}{dx} = \frac{1}{\cos y}
  \]
  Bây giờ, $\cos y \ge 0$ vì $-\pi/2 \le y \le \pi/2$, do đó
  \[
    \cos y = \sqrt{1 - \sin^2 y} = \sqrt{1 - x^2} \qquad {\color{stewartcyan}(\cos^2 y + \sin^2 y = 1)}
  \]
  \noindent
  \makebox[0pt][l]{Do đó}%
  \hfill $\displaystyle\frac{dy}{dx} = \frac{1}{\cos y} = \frac{1}{\sqrt{1 - x^2}}$\hfill\null

  \vspace{5pt}
  \begin{center}
    \fcolorbox{stewartred}{white}{%
      \hspace{1.8em}$\displaystyle \frac{d}{dx}(\sin^{-1} x) = \frac{1}{\sqrt{1 - x^2}}$\hspace{1.8em}%
    }
  \end{center}
\end{minipage}

\end{document}